% 图 81.1 —— 由 `checks/emit_hand.py` 自动拼出（probe 精测 + prims2 图元分解）
%%
% ⚠ 这是**稿子不是成品**：跑 `checks/checkhand.py 81.1` 看指标 + 看叠合图。
%%
%% ⚠ 待核：
%   · 本图 probe 报出阴影族 1 个 —— **本脚本不生成阴影**（要照原书手写）；prims2 的 strokes 里可能含阴影线，核对时留意别画重
%   · 可疑字形 5 项（空心圈 ⊙ 常被认成字母 o）—— 见文件末
%%
\documentclass[12pt]{standalone}
\usepackage{fontspec}
\setsansfont{Arial}
\newfontfamily\mfallback{DejaVu Sans}
\usepackage{tikz}
\usetikzlibrary{arrows.meta}
\begin{document}
\begin{tikzpicture}[x=1pt, y=-1pt, line width=5.0pt, line cap=round, line join=round]

  %% ════ 填充区（prims2 的 fills）════
  \fill (78.0,466.0) -- (78.0,470.0) -- (91.0,483.0) -- (93.0,481.0) -- (81.0,466.0) -- cycle;

  %% ════ 直线段（prims2 的 strokes）════
  \draw[line width=5.0pt] (198.9,26.0) -- (187.0,31.0);
  \draw[line width=5.0pt] (187.0,31.0) -- (168.0,41.0);
  \draw[line width=4.0pt] (168.0,41.0) -- (136.0,63.0);
  \draw[line width=7.0pt] (281.0,85.0) -- (281.0,76.0);
  \draw[line width=5.0pt] (103.0,449.0) -- (109.0,467.0);
  \draw[line width=5.7pt] (277.0,132.0) -- (282.0,130.0);
  \draw[line width=5.0pt] (101.0,415.0) -- (101.0,432.0);
  \draw[line width=5.0pt] (87.0,249.0) -- (78.0,265.0);
  \draw[line width=5.0pt] (35.0,324.0) -- (27.0,330.0);
  \draw[line width=7.7pt] (277.0,389.0) -- (284.0,389.0);
  \draw[line width=3.0pt] (285.0,229.0) -- (285.0,232.0);
  \draw[line width=4.0pt] (546.0,353.0) -- (447.0,298.0);
  \draw[line width=7.0pt] (796.0,296.0) -- (790.0,293.0);
  \draw[line width=6.3pt] (642.0,519.0) -- (648.0,518.0);
  \draw[line width=5.0pt] (59.0,308.0) -- (55.0,326.0);
  \draw[line width=5.0pt] (107.0,225.0) -- (96.0,237.0);
  \draw[line width=7.0pt] (117.0,372.0) -- (126.0,357.0);
  \draw[line width=8.1pt] (100.0,223.0) -- (107.0,224.0);
  \draw[line width=7.0pt] (108.0,224.0) -- (124.0,210.0);
  \draw[line width=6.5pt] (284.0,130.0) -- (289.0,134.0);
  \draw[line width=7.3pt] (548.0,353.0) -- (547.0,346.0);
  \draw[line width=5.0pt] (112.0,179.0) -- (126.0,191.0);
  \draw[line width=6.0pt] (117.0,373.0) -- (118.0,385.0);
  \draw[line width=6.0pt] (32.0,338.0) -- (27.0,330.0);
  \draw[line width=5.0pt] (80.0,466.0) -- (90.0,480.0);
  \draw[line width=5.7pt] (790.0,484.0) -- (785.0,482.0);
  \draw[line width=5.0pt] (790.0,484.0) -- (784.3,490.7);
  \draw[line width=5.0pt] (105.0,166.0) -- (100.0,149.0);
  \draw[line width=11.0pt] (101.0,96.0) -- (111.0,95.0);
  \draw[line width=6.5pt] (293.0,292.0) -- (288.0,288.0);
  \draw[line width=4.0pt] (649.0,517.0) -- (663.9,478.1);
  \draw[line width=5.0pt] (52.0,406.0) -- (58.0,425.0);
  \draw[line width=4.0pt] (50.0,394.0) -- (50.0,373.0);
  \draw[line width=5.0pt] (140.0,197.0) -- (167.0,180.0);
  \draw[line width=5.0pt] (167.0,180.0) -- (186.9,170.0);
  \draw[line width=5.0pt] (186.9,170.0) -- (210.7,162.0);
  \draw[line width=4.0pt] (51.0,343.0) -- (49.0,361.0);
  \draw[line width=5.0pt] (696.1,268.0) -- (671.0,294.0);
  \draw[line width=5.0pt] (110.0,97.0) -- (106.0,102.0);
  \draw[line width=6.0pt] (25.0,330.0) -- (21.0,336.0);
  \draw[line width=5.0pt] (120.0,497.0) -- (102.0,490.0);
  \draw[line width=5.0pt] (255.0,353.0) -- (285.0,290.0);
  \draw[line width=5.0pt] (26.0,318.0) -- (25.0,329.0);
  \draw[line width=7.0pt] (112.0,97.0) -- (115.0,104.0);
  \draw[line width=5.0pt] (152.2,325.8) -- (143.0,344.0);
  \draw[line width=4.0pt] (101.0,116.0) -- (99.0,134.0);
  \draw[line width=7.5pt] (106.0,374.0) -- (116.0,373.0);
  \draw[line width=4.0pt] (276.0,390.0) -- (276.0,394.0);
  \draw[line width=4.0pt] (284.0,427.9) -- (286.0,415.7);
  \draw[line width=4.0pt] (71.0,454.0) -- (64.0,440.0);
  \draw[line width=6.0pt] (540.0,357.0) -- (546.0,354.0);
  \draw[line width=5.3pt] (793.0,484.0) -- (794.0,489.0);
  \draw[line width=8.7pt] (110.0,233.0) -- (108.0,225.0);
  \draw[line width=7.0pt] (112.0,95.0) -- (122.0,77.0);
  \draw[line width=7.0pt] (560.0,361.0) -- (548.0,354.0);
  \draw[line width=4.5pt] (71.0,278.0) -- (63.0,296.0);
  \draw[line width=4.5pt] (80.0,399.0) -- (75.0,412.0);
  \draw[line width=4.0pt] (276.0,154.0) -- (283.0,131.0);

  %% ════ 曲线（prims2 的 curves —— 三段式，坐标同为 (y,x)）════
  \draw[line width=5.0pt] (281.0,74.0) .. controls (290.0,30.6) and (232.4,4.7) .. (198.9,26.0);
  \draw[line width=5.0pt] (281.0,85.0) .. controls (271.6,95.0) and (252.6,89.5) .. (252.0,74.8);
  \draw[line width=5.0pt] (252.0,74.8) .. controls (255.8,60.7) and (272.7,68.5) .. (280.0,75.0);
  \draw[line width=4.0pt] (281.0,85.0) .. controls (288.1,97.7) and (285.3,114.8) .. (283.0,129.0);
  \draw[line width=4.0pt] (284.0,228.0) .. controls (275.9,219.8) and (256.9,219.7) .. (260.0,235.5);
  \draw[line width=5.0pt] (260.0,235.5) .. controls (264.6,247.9) and (278.6,241.6) .. (284.0,233.0);
  \draw[line width=4.0pt] (651.0,518.0) .. controls (685.0,489.9) and (714.4,458.4) .. (742.0,424.0);
  \draw[line width=7.0pt] (641.0,520.0) .. controls (638.4,526.8) and (647.2,528.2) .. (650.0,523.0);
  \draw[line width=4.0pt] (285.0,233.0) .. controls (298.8,248.1) and (293.1,269.8) .. (288.0,287.0);
  \draw[line width=4.0pt] (789.0,294.0) .. controls (794.0,313.6) and (783.4,336.0) .. (766.0,347.0);
  \draw[line width=6.0pt] (81.0,399.0) .. controls (89.6,396.0) and (89.4,406.3) .. (86.0,411.0);
  \draw[line width=5.0pt] (784.3,490.7) .. controls (745.4,515.0) and (697.9,521.4) .. (652.0,522.0);
  \draw[line width=4.0pt] (663.9,478.1) .. controls (680.9,429.6) and (728.6,384.9) .. (783.9,395.0);
  \draw[line width=4.0pt] (783.9,395.0) .. controls (821.1,407.5) and (818.7,460.3) .. (793.0,483.0);
  \draw[line width=5.0pt] (210.7,162.0) .. controls (249.5,146.7) and (308.7,186.0) .. (285.0,228.0);
  \draw[line width=5.0pt] (789.0,292.0) .. controls (787.2,278.1) and (783.8,263.4) .. (771.0,255.0);
  \draw[line width=4.0pt] (771.0,255.0) .. controls (761.5,248.7) and (749.9,246.8) .. (737.9,248.1);
  \draw[line width=4.0pt] (737.9,248.1) .. controls (723.0,251.7) and (708.8,259.0) .. (696.1,268.0);
  \draw[line width=5.0pt] (671.0,294.0) .. controls (632.2,359.9) and (619.6,444.6) .. (641.0,519.0);
  \draw[line width=5.0pt] (262.0,187.0) .. controls (239.8,240.8) and (189.0,279.7) .. (152.2,325.8);
  \draw[line width=5.0pt] (143.0,344.0) .. controls (190.2,348.1) and (242.0,350.8) .. (275.0,389.0);
  \draw[line width=4.0pt] (759.0,383.0) .. controls (766.0,374.3) and (767.7,361.8) .. (765.0,351.0);
  \draw[line width=5.0pt] (137.0,490.0) .. controls (152.8,439.9) and (216.6,430.1) .. (238.0,384.0);
  \draw[line width=4.0pt] (137.0,490.0) .. controls (128.2,491.1) and (122.7,483.7) .. (117.0,478.0);
  \fill (140.0,485.0) -- (134.0,495.0) -- (122.1,481.0) -- cycle;
  \draw[line width=5.0pt] (137.0,490.0) .. controls (190.4,488.7) and (255.7,480.7) .. (284.0,427.9);
  \draw[line width=4.0pt] (286.0,415.7) .. controls (284.4,408.5) and (284.5,399.1) .. (277.0,395.0);
  \draw[line width=4.0pt] (765.0,347.0) .. controls (761.9,339.7) and (754.1,327.1) .. (745.4,334.6);
  \draw[line width=5.0pt] (745.4,334.6) .. controls (738.7,346.7) and (753.8,355.1) .. (764.0,351.0);

  %% ════ 实心点 ════
  \fill (139.6,60.4) circle (11.2);
  \fill (139.6,198.2) circle (7.0);
  \fill (785.0,296.5) circle (4.5);
  \fill (142.5,345.1) circle (7.1);
  % （略）(134.5,493.0) r=6.8 落在箭头头上 —— 已由箭头头代替

  %% ════ 标签（probe 直接给的 \node 行）════
  %% ⚠ 全部补了 fill=white：文字在最上层，否则与曲线交叉干扰
     \node[anchor=center,inner sep=0pt,fill=white,font=\fontsize{29.5}{36.9}\selectfont] at (152.7,89.5) {\textsf{\textbf{\textit{o}}}};   % box(144,80,16,16)
     \node[anchor=center,inner sep=0pt,fill=white,font=\fontsize{23.0}{28.7}\selectfont] at (165.9,99.1) {\textsf{\textbf{3}}};   % box(160,90,11,17)
     \node[anchor=center,inner sep=0pt,fill=white,font=\fontsize{34.4}{43.0}\selectfont] at (354.4,126.4) {\textsf{\textbf{6}}};   % box(346,112,15,27)
     \node[anchor=center,inner sep=0pt,fill=white,font=\fontsize{30.0}{37.5}\selectfont] at (314.8,128.2) {\textsf{\textbf{\textit{b}}}};   % box(305,116,17,22)
     \node[anchor=center,inner sep=0pt,fill=white,font=\fontsize{29.1}{36.4}\selectfont] at (75.2,130.2) {\textsf{\textbf{\textit{b}}}};   % box(66,118,16,22)
     \node[anchor=center,inner sep=0pt,fill=white,font=\fontsize{17.2}{21.5}\selectfont] at (333.7,128.1) {\textsf{\textbf{\textit{O}}}};   % box(327,121,12,13)
     \node[anchor=center,inner sep=0pt,fill=white,font=\fontsize{31.4}{39.3}\selectfont] at (154.8,223.1) {\textsf{\textbf{\textit{e}}}};   % box(145,213,16,17)
     \node[anchor=center,inner sep=0pt,fill=white,font=\fontsize{23.6}{29.6}\selectfont] at (166.1,232.9) {\textsf{\textbf{2}}};   % box(160,224,11,17)
     \node[anchor=center,inner sep=0pt,fill=white,font=\fontsize{33.8}{42.2}\selectfont] at (316.3,295.8) {\textsf{\textbf{6}}};   % box(308,282,15,26)
     \node[anchor=center,inner sep=0pt,fill=white,font=\fontsize{29.8}{37.3}\selectfont] at (514.0,301.5) {\textsf{\textbf{\textit{P}}}};   % box(503,289,18,23)
     \node[anchor=center,inner sep=0pt,fill=white,font=\fontsize{32.2}{40.2}\selectfont] at (819.4,317.0) {\textsf{\textbf{\textit{B}}}};   % box(810,300,15,33)
     \node[anchor=center,inner sep=0pt,fill=white,font=\fontsize{31.3}{39.2}\selectfont] at (152.3,374.6) {\textsf{\textbf{\textit{e}}}};   % box(143,364,15,18)
     \node[anchor=center,inner sep=0pt,fill=white,font=\fontsize{21.5}{26.9}\selectfont] at (165.0,384.9) {\textsf{\textbf{1}}};   % box(159,377,7,15)
     \node[anchor=center,inner sep=0pt,fill=white,font=\fontsize{29.6}{37.0}\selectfont] at (317.0,395.0) {\textsf{\textbf{\textit{Y}}}};   % box(308,382,17,24)
     \node[anchor=center,inner sep=0pt,fill=white,font=\fontsize{27.8}{34.8}\selectfont] at (109.0,396.0) {\textsf{\textbf{\textit{f}}}};   % box(103,385,11,21)
     \node[anchor=center,inner sep=0pt,fill=white,font=\fontsize{32.6}{40.8}\selectfont] at (833.0,461.5) {\textsf{\textbf{\textit{α}}}};   % box(824,452,19,18)
     \node[anchor=center,inner sep=0pt,fill=white,font=\fontsize{31.4}{39.3}\selectfont] at (141.8,520.1) {\textsf{\textbf{\textit{e}}}};   % box(132,510,16,17)
     \node[anchor=center,inner sep=0pt,fill=white,font=\fontsize{23.8}{29.8}\selectfont] at (154.5,529.7) {\textsf{\textbf{0}}};   % box(147,521,11,17)

  % 钉死画布：否则超出的图元/控制点会把页面撑大，1:1 回比就失准
  \pgfresetboundingbox
  \path[use as bounding box] (0,0) rectangle (856,546);
\end{tikzpicture}
\end{document}

% ⚠ 待核清单（probe 拿不准，人眼过一遍）：
%   · 可疑字形：\node[anchor=center,inner sep=0pt,font=\fontsize{29.5}{36.9}\selectfont] at (152.7,89.5) {\textsf{\textbf{\tex
%   · 可疑字形：\node[anchor=center,inner sep=0pt,font=\fontsize{17.2}{21.5}\selectfont] at (333.7,128.1) {\textsf{\textbf{\te
%   · 未识别/跳过：⑤ 跳过/未识别的字形
%   · 未识别/跳过：box(111,179,18,16)
%   · 未识别/跳过：box(101,488,21,11)
